\documentclass[12pt,a4paper]{article}
\usepackage[utf8]{inputenc}
\usepackage[T5]{fontenc} % Sửa lỗi font chữ tiếng Việt bị mờ/xấu (bitmapped)
\usepackage[vietnamese]{babel}
\usepackage{mathptmx} % Sử dụng font Times New Roman chuẩn học thuật
\usepackage{setspace}
\onehalfspacing % Giãn dòng 1.5 chuẩn KLTN
\usepackage{amsmath}
\usepackage{graphicx}
\usepackage{booktabs}
\usepackage{float}
\usepackage{tikz}
\usetikzlibrary{arrows.meta,calc,positioning}
\usepackage{hyperref}
\usepackage{listings}
\usepackage{geometry}
\geometry{
 a4paper,
 total={170mm,257mm},
 left=30mm, % Lề trái 3cm chuẩn đóng bìa
 right=20mm,
 top=20mm,
 bottom=20mm,
}

% Màu và kiểu dùng chung cho các sơ đồ kiến trúc phần cứng.
\definecolor{u250blue}{HTML}{D9EAF7}
\definecolor{u250navy}{HTML}{28587B}
\definecolor{u250green}{HTML}{DDF2E1}
\definecolor{u250orange}{HTML}{FCE8CE}
\definecolor{u250purple}{HTML}{E9E0F5}
\definecolor{u250gray}{HTML}{EEF1F4}
\definecolor{u250red}{HTML}{F6D7D7}

\tikzset{
    slrbox/.style={draw=u250navy, very thick, rounded corners=3pt,
        fill=u250blue!35, minimum width=3.8cm, minimum height=5.25cm},
    ddrbox/.style={draw=u250navy, thick, rounded corners=2pt,
        fill=u250gray, minimum width=2.55cm, minimum height=0.72cm,
        align=center, font=\small\bfseries},
    pebox/.style={draw=u250navy, thick, rounded corners=2pt,
        fill=u250green, minimum width=3.05cm, minimum height=1.35cm,
        align=center, font=\small},
    localbox/.style={draw=u250navy, thick, rounded corners=2pt,
        fill=white, minimum width=3.05cm, minimum height=1.15cm,
        align=center, font=\scriptsize},
    reducebox/.style={draw=u250navy, thick, rounded corners=2pt,
        fill=u250purple, minimum width=3.25cm, minimum height=0.85cm,
        align=center, font=\small},
    hwbox/.style={draw=u250navy, thick, rounded corners=2pt,
        fill=u250gray, minimum height=0.9cm, align=center, font=\small},
    dspbox/.style={draw=u250navy, very thick, rounded corners=2pt,
        fill=u250orange, minimum width=2.3cm, minimum height=1.25cm,
        align=center, font=\small\bfseries},
    streambox/.style={draw=u250navy, thick, rounded corners=2pt,
        fill=u250green, minimum width=1.8cm, minimum height=1.05cm,
        align=center, font=\scriptsize},
    dataline/.style={-{Latex[length=2.1mm,width=1.5mm]}, very thick,
        draw=u250navy},
    cmdline/.style={-{Latex[length=1.8mm,width=1.2mm]}, thick, dashed,
        draw=black!65},
    bothline/.style={{Latex[length=1.8mm,width=1.2mm]}-{Latex[length=1.8mm,width=1.2mm]},
        very thick, draw=u250navy}
}

\title{\textbf{Báo Cáo Kiến Trúc Phần Cứng}\\ \Large Bộ Giải Mã Llama-2 7B INT4 trên FPGA Alveo U250}
\author{Khóa Luận Tốt Nghiệp}
\date{\today}

\begin{document}

\maketitle
\tableofcontents
\newpage

\section{Giới thiệu chung}
Báo cáo này mô tả chi tiết kiến trúc phần cứng của hệ thống giải mã (decoder) token Llama-2 7B được lượng tử hóa xuống chuẩn INT4, triển khai trên FPGA Xilinx Alveo U250 (\texttt{xcu250-figd2104-2L-e}) sử dụng Vitis HLS 2023.2. 

Kiến trúc được phân mảnh (sharded) theo cột (input column) và phân bố đều trên 4 Super Logic Regions (SLR) của card U250, tương ứng với 4 Processing Elements (PE) hoạt động độc lập với 4 bộ nhớ DDR (DDR0 đến DDR3). 

Nguyên tắc thiết kế bắt buộc:
\begin{itemize}
    \item Mỗi PE/SLR/DDR là một miền dữ liệu độc lập. Không PE nào truy cập chéo bộ nhớ tham số (model), RoPE, residual hay KV cache của PE khác.
    \item Chỉ các luồng dữ liệu trung gian (partial output 128-bit) và lệnh điều khiển (scalar command/reduction) được phép đi qua ranh giới giữa các SLR để tránh hiện tượng nghẽn cổ chai vật lý.
\end{itemize}

\section{Kiến trúc Hệ thống (Top-Level Module)}
Kernel chính điều khiển toàn bộ luồng dữ liệu là \texttt{int4\_decoder\_token\_controller}. Khi được gọi, kernel này nhận một hidden state đã được nhúng (embedding), sau đó chạy tuần tự qua 32 lớp decoder, cập nhật KV cache cục bộ và xuất ra logits.

\subsection{Giao tiếp và Bộ nhớ}
Mỗi PE quản lý một cổng giao tiếp AXI với bộ nhớ DDR cục bộ của nó:
\begin{itemize}
    \item \textbf{Model Bank (\texttt{gmem0-3}):} Lưu trữ trọng số mô hình (weights), scales và các hệ số chuẩn hóa. Tổng cộng khoảng $13.1$ triệu từ 512-bit trên mỗi DDR.
    \item \textbf{RoPE LUT:} Bảng tra cứu cho Rotary Positional Embedding.
    \item \textbf{KV Cache:} Lưu trữ Key và Value cục bộ phục vụ cơ chế Attention.
    \item \textbf{Residual \& Logits:} Dữ liệu vào/ra chính của pipeline.
\end{itemize}

Dữ liệu tạm thời được đưa vào bộ nhớ nội bộ của FPGA:
\begin{itemize}
    \item \textbf{URAM:} Dùng để lưu trữ bộ nhớ đệm (cache) cho Scale và Norm Gamma nhằm đảm bảo tốc độ truy xuất ngẫu nhiên nhanh.
    \item \textbf{BRAM:} Làm bộ nhớ nháp (scratchpad) lưu các giá trị trung gian như Q, K, V, Gate, và các luồng Residual.
\end{itemize}

\section{Kiến trúc Xử lý Nội bộ và Các Sub-Module}
Hệ thống sử dụng cơ chế \textbf{Time-Multiplexing} (Tái sử dụng phần cứng). Thay vì nhân bản phần cứng cho từng lớp hoặc từng phép tính, toàn bộ pipeline dùng chung một khối nhân ma trận động.

\subsection{Động cơ Nhân ma trận Tensor-Parallel (\texttt{int4\_sharded\_linear\_4pe})}
Khối Linear là thành phần lõi của thiết kế, được sử dụng lại cho toàn bộ các phép chiếu (projections): Q, K, V, O, Gate, Up, Down và Logits. 

Bốn PE độc lập tự động lấy dữ liệu từ BRAM cục bộ và DDR, tính toán các tổng thành phần (partial sum). Các tổng này (128-bit) được chuyển chéo qua các SLR qua một mạng lưới Reduction (Cộng gộp):
\begin{itemize}
    \item \textbf{SLR1 (Pair 01):} Cộng gộp kết quả của PE0 và PE1.
    \item \textbf{SLR2 (Pair 23):} Cộng gộp kết quả của PE2 và PE3.
\end{itemize}
Sự kết hợp chéo giữa Pair01 và Pair23 sẽ tạo ra output hoàn chỉnh, sau đó được phát ngược lại vào BRAM của mỗi SLR.

\subsection{Kỹ thuật Đóng gói DSP48 (DSP Packing)}
Điểm nhấn kỹ thuật giúp tiết kiệm tài nguyên là việc nhồi \textbf{hai phép nhân \texttt{INT4 x INT15}} vào cùng một khối DSP48. 

Hai trọng số INT4 của hai hàng ngõ ra chia sẻ chung một giá trị kích hoạt (activation INT15). Trọng số được đóng gói vào toán hạng 27-bit, cách nhau 23-bit:
\begin{equation}
    packed\_weight = (w_{high} \ll 23) + w_{low}
\end{equation}
\begin{equation}
    packed\_weight \times activation = (w_{high} \times activation \ll 23) + (w_{low} \times activation)
\end{equation}
Khoảng cách 23-bit đủ lớn để chứa tổng có dấu của một group 32 phần tử (tương ứng với một SIMD lane) mà không bị tràn (overflow). Điều này cho phép một khối DSP48 xử lý đồng thời 2 phép tính, nâng cao thông lượng lên 128 phép nhân scalar mỗi chu kỳ cho một PE, mà chỉ chiếm 7\% tài nguyên DSP của chip.

\subsection{Cơ chế Nhân Ma Trận với Vector (GEMV) trên Phần Cứng}
Mạng nơ-ron thực chất là các phép nhân ma trận trọng số $W (M \times N)$ với vector kích hoạt $x (N \times 1)$. Trên FPGA, thay vì dùng các vòng lặp phần mềm chậm chạp, phép toán GEMV (General Matrix-Vector Multiplication) được chuyển hóa thành mạch logic không gian (Spatial Computing) cực kỳ song song:

\begin{itemize}
    \item \textbf{Cấu trúc dữ liệu siêu rộng (Wide Memory Access):} Trọng số INT4 không được lưu trữ theo mảng 1D thông thường, mà được tổ chức theo bố cục \textit{Transposed G32}. Mỗi từ nhớ (word) đọc từ DDR có kích thước khổng lồ 512-bit, chứa dữ liệu của 4 hàng và 32 cột ($4 \times 32 \times 4\text{-bit} = 512\text{-bit}$). Nghĩa là mỗi nhịp clock, FPGA nạp vào trọn vẹn 128 trọng số.
    \item \textbf{Mở phẳng vòng lặp (Loop Unrolling):} Vòng lặp tính toán dọc theo chiều cột (K-dimension) được mở phẳng hoàn toàn (\texttt{\#pragma HLS UNROLL}). Phần cứng sinh ra 32 làn (lane) nhân-cộng hoạt động hoàn toàn song song.
    \item \textbf{Triển khai trên Mạch:} Trong một chu kỳ xung nhịp (clock cycle), 1 vector kích hoạt (gồm 32 phần tử) và 1 word trọng số 512-bit được đưa vào 64 khối DSP48 (mỗi khối xử lý 2 hàng nhờ DSP Packing). FPGA thực hiện đồng thời $32 \times 4 = 128$ phép nhân-cộng (MAC).
    \item \textbf{Cây cộng (Adder Tree):} Kết quả của 32 phép nhân không dùng vòng lặp để cộng dồn, mà đi qua một cấu trúc mạch "Cây cộng" (Adder Tree) để ra ngay tổng tích lũy của 4 hàng trong cùng một chu kỳ, rồi cập nhật vào BRAM. Quá trình này duy trì thông lượng 1 vòng lặp / chu kỳ (Initiation Interval II=1).
\end{itemize}

\subsection{Kỹ thuật Che Khuất Độ Trễ và Gối Đầu (Overlapping bằng DATAFLOW)}
Rào cản lớn nhất trong suy luận LLM là thời gian chờ dữ liệu từ bộ nhớ DDR (Memory Latency), thường được gọi là "Memory-bound". Để phần cứng luôn đạt 100% công suất (không bị stall), hệ thống áp dụng kỹ thuật \textbf{Pipelining mức Task} thông qua chỉ thị \texttt{\#pragma HLS DATAFLOW}.

Quá trình nhân một khối ma trận (Tile) được cắt thành 3 tiến trình (Task) hoạt động độc lập và nối với nhau bằng các ống nước (hls::stream FIFO):
\begin{enumerate}
    \item \textbf{Tiến trình Load (Đọc):} Module \texttt{stream\_weight\_tiles} liên tục gửi lệnh AXI Burst để hút dữ liệu từ DDR, đẩy vào một FIFO đệm (\texttt{weight\_stream}) có độ sâu 512 words (đủ chứa 2 Tile ma trận).
    \item \textbf{Tiến trình Compute (Tính toán):} Module lõi \texttt{compute\_streamed\_mvm} không quan tâm đến DDR. Nó chỉ lấy trọng số từ FIFO và kích hoạt từ BRAM cục bộ, sau đó vắt kiệt sức mạnh của DSP ở tốc độ II=1.
    \item \textbf{Tiến trình Store (Ghi/Reduce):} Lấy kết quả đã cộng (partial sums) đẩy ra mạng Reduction liên SLR.
\end{enumerate}

\textbf{Hiệu ứng Overlap (Gối đầu):} Nhờ DATAFLOW và FIFO đủ lớn (Elastic Buffer), 3 tiến trình này chạy song song chồng lấn lên nhau. Tại bất kỳ thời điểm nào, phần cứng cũng đang: \textbf{Load dữ liệu cho Tile $n+1$}, đồng thời \textbf{Compute Tile $n$}, và \textbf{Store kết quả của Tile $n-1$}. Độ trễ biến động của DDR (Burst Jitter) hoàn toàn bị "nuốt chửng" bởi các FIFO, giúp lõi tính toán không bao giờ phải chờ đợi (zero-idle).

\subsection{Các Khối Xử lý Khác}
\begin{itemize}
    \item \textbf{\texttt{int4\_rmsnorm\_quantize\_shards}:} Thực hiện chuẩn hóa RMS. Đây là khối duy nhất truyền dữ liệu liên SLR gồm 1 số FP32 (tổng bình phương) và 1 nghịch đảo căn bậc hai. Sau khi chuẩn hóa, kết quả được lượng tử hóa sang dạng INT4.
    \item \textbf{\texttt{int4\_swiftkv\_attention\_pe}:} Tính toán Attention cục bộ trên từng SLR. Do dữ liệu K, V đã được chia nhỏ theo head, module này chạy hoàn toàn độc lập, sử dụng thuật toán online-softmax.
    \item \textbf{\texttt{int4\_swiglu\_quantize\_shards}:} Tính hàm kích hoạt $\text{Swish}(\text{Gate}) \times \text{Up}$ cho khối Feed-Forward (FFN), sau đó lượng tử hóa kết quả.
\end{itemize}

\section{Luồng Dữ Liệu (Dataflow)}
Để khắc phục hạn chế về đồng bộ của kiến trúc HLS DATAFLOW phân tán trên nhiều SLR, hệ thống sử dụng một chuỗi Token chuyền qua các PE (Token Control Network) thay vì sử dụng máy trạng thái trung tâm.

Quy trình hoạt động cho 1 Layer (Lặp lại 32 lần):
\begin{enumerate}
    \item \textbf{RMS(attn)} $\rightarrow$ Sinh ra Activation lượng tử hóa.
    \item Khối \textbf{Linear} lần lượt tính \textbf{Q, K, V}.
    \item Khối \textbf{Attention (SwiftKV + RoPE)} tính toán ngữ cảnh, cập nhật bộ nhớ DDR.
    \item Khối \textbf{Linear} chiếu kết quả tạo \textbf{O}, sau đó cộng \textbf{Residual}.
    \item \textbf{RMS(ffn)} $\rightarrow$ Khối \textbf{Linear} tính \textbf{Gate, Up}.
    \item Khối \textbf{SwiGLU} kích hoạt phi tuyến.
    \item Khối \textbf{Linear} chiếu \textbf{Down}, sau đó cộng \textbf{Residual}.
\end{enumerate}

\section{Sử dụng Tài nguyên và Hiệu năng}
Dựa trên báo cáo tổng hợp C-Synthesis (Vitis HLS 2023.2):
\begin{table}[H]
    \centering
    \begin{tabular}{l l r}
    \toprule
    \textbf{Chỉ số} & \textbf{Chi tiết} & \textbf{Kết quả} \\
    \midrule
    Mục tiêu xung nhịp & Target Clock & $300$ MHz ($3.333$ ns) \\
    Initiation Interval & Linear MAC II & 1 \\
    Đóng gói MAC & Lượng tính toán / DSP48 & 2 phép nhân \texttt{INT4xINT15} \\
    Sử dụng BRAM18K & Bộ nhớ nội bộ (Scratchpad) & $1308 / 5376$ ($24\%$) \\
    Sử dụng DSP & Khối tính toán & $916 / 12288$ ($7\%$) \\
    Sử dụng LUT & Logic tùy biến & $387973 / 1728000$ ($22\%$) \\
    Sử dụng URAM & Bộ nhớ cache (Scales/Norms) & $160 / 1280$ ($12\%$) \\
    \bottomrule
    \end{tabular}
    \caption{Tóm tắt tài nguyên phần cứng ước tính}
\end{table}
Với tài nguyên sử dụng tối ưu, mạng lưới floorplan được thiết kế cứng (\texttt{hard floorplan}) thông qua file cấu hình TCL để đảm bảo đạt tiêu chuẩn timing khắt khe ở mức $300$ MHz mà không xảy ra lỗi định tuyến liên SLR.

\clearpage
\section{Sơ đồ kiến trúc phần cứng}

Hai sơ đồ dưới đây được rút trực tiếp từ cấu trúc module và ràng buộc kết nối
trong source.  Mỗi PE sở hữu một miền DDR/SLR riêng, gồm model shard, activation
scratchpad và KV cache cục bộ.  Luồng điều khiển xuyên SLR chỉ mang lệnh scalar;
đường dữ liệu xuyên SLR có bề rộng lớn duy nhất là các packet partial 128 bit
tại tầng reduction.

\begin{figure}[H]
\centering
\resizebox{\textwidth}{!}{%
\begin{tikzpicture}[x=1cm,y=1cm]
    % Bốn DDR nằm ngoài FPGA và chỉ nối với SLR tương ứng.
    \node[ddrbox] (ddr0) at (0,7.15) {DDR0 / gmem0\\AXI 512 bit};
    \node[ddrbox] (ddr1) at (4.5,7.15) {DDR1 / gmem1\\AXI 512 bit};
    \node[ddrbox] (ddr2) at (9.0,7.15) {DDR2 / gmem2\\AXI 512 bit};
    \node[ddrbox] (ddr3) at (13.5,7.15) {DDR3 / gmem3\\AXI 512 bit};

    % Nền SLR được vẽ trước các khối con.
    \node[slrbox] (slr0) at (0,3.85) {};
    \node[slrbox] (slr1) at (4.5,3.85) {};
    \node[slrbox] (slr2) at (9.0,3.85) {};
    \node[slrbox] (slr3) at (13.5,3.85) {};

    \node[font=\bfseries\small, text=u250navy] at (0,6.15) {SLR0};
    \node[font=\bfseries\small, text=u250navy] at (4.5,6.15) {SLR1};
    \node[font=\bfseries\small, text=u250navy] at (9.0,6.15) {SLR2};
    \node[font=\bfseries\small, text=u250navy] at (13.5,6.15) {SLR3};

    \node[pebox] (pe0) at (0,4.75) {\textbf{Linear PE0}\\RMSNorm, RoPE, SwiftKV\\SwiGLU, residual};
    \node[pebox] (pe1) at (4.5,4.75) {\textbf{Linear PE1}\\RMSNorm, RoPE, SwiftKV\\SwiGLU, residual};
    \node[pebox] (pe2) at (9.0,4.75) {\textbf{Linear PE2}\\RMSNorm, RoPE, SwiftKV\\SwiGLU, residual};
    \node[pebox] (pe3) at (13.5,4.75) {\textbf{Linear PE3}\\RMSNorm, RoPE, SwiftKV\\SwiGLU, residual};

    \node[localbox] (mem0) at (0,2.85) {BRAM/URAM cục bộ\\model shard, scale, activation\\KV cache, residual, logits};
    \node[localbox] (mem1) at (4.5,2.85) {BRAM/URAM cục bộ\\model shard, scale, activation\\KV cache, residual, logits};
    \node[localbox] (mem2) at (9.0,2.85) {BRAM/URAM cục bộ\\model shard, scale, activation\\KV cache, residual, logits};
    \node[localbox] (mem3) at (13.5,2.85) {BRAM/URAM cục bộ\\model shard, scale, activation\\KV cache, residual, logits};

    \draw[dataline] (ddr0) -- node[right, font=\scriptsize] {weight} (pe0);
    \draw[dataline] (ddr1) -- node[right, font=\scriptsize] {weight} (pe1);
    \draw[dataline] (ddr2) -- node[right, font=\scriptsize] {weight} (pe2);
    \draw[dataline] (ddr3) -- node[right, font=\scriptsize] {weight} (pe3);
    \draw[bothline] (pe0) -- (mem0);
    \draw[bothline] (pe1) -- (mem1);
    \draw[bothline] (pe2) -- (mem2);
    \draw[bothline] (pe3) -- (mem3);

    % Command chain rất hẹp, tránh fanout điều khiển xuyên SLR.
    \draw[cmdline] (-1.45,1.85) -- (1.45,1.85);
    \draw[cmdline] (1.65,1.85) -- (6.15,1.85);
    \draw[cmdline] (6.35,1.85) -- (10.65,1.85);
    \draw[cmdline] (10.85,1.85) -- (15.05,1.85);
    \node[font=\scriptsize, fill=white, inner sep=2pt] at (6.75,2.15)
        {chuỗi lệnh scalar: PE0 $\rightarrow$ PE1 $\rightarrow$ PE2 $\rightarrow$ PE3};

    % Mạng cộng logic: chỉ các gói partial đi qua biên SLR.
    \node[reducebox] (r01) at (2.25,0.45) {pair reduction 01\\FP32 partial, 128 bit};
    \node[reducebox] (r23) at (11.25,0.45) {pair reduction 23\\FP32 partial, 128 bit};
    \draw[dataline] (mem0.south) -- (r01.north west);
    \draw[dataline] (mem1.south) -- (r01.north east);
    \draw[dataline] (mem2.south) -- (r23.north west);
    \draw[dataline] (mem3.south) -- (r23.north east);
    \draw[bothline] (r01) -- node[above, font=\scriptsize, fill=white]
        {mỗi chiều một packet 128 bit} (r23);

    \node[font=\scriptsize, align=center] at (6.75,-0.55)
        {Cộng đủ 4 shard rồi ghi/broadcast kết quả về projection scratch cục bộ;\\
         scheduler tái sử dụng cùng datapath cho Q, K, V, O, Gate, Up, Down và Logits.};
\end{tikzpicture}%
}
\caption{Kiến trúc tổng quan của bộ tăng tốc Llama 2-7B trên Alveo U250.}
\label{fig:u250-overview}
\end{figure}

Hình~\ref{fig:w4a15-dataflow} ghép hai mức nhìn của cùng datapath. Phần (a)
mô tả phép đóng gói signed bên trong một DSP48E2. Phần (b) cho thấy FIFO BRAM
tách AXI burst khỏi pipeline MAC, nhờ đó tile $n+1$ được nạp trong khi tile
$n$ đang tính và packet trước đó đang được reduction/ghi ra.

\begin{figure}[H]
\centering
\resizebox{\textwidth}{!}{%
\begin{tikzpicture}[x=1cm,y=1cm]
    \node[font=\bfseries, anchor=west] at (-0.1,6.35)
        {(a) Signed W4A15 packing trong một DSP48E2};

    \node[hwbox, fill=u250red, minimum width=1.55cm] (wh) at (0.8,5.5)
        {$w_{\mathrm{hi}}$\\signed W4};
    \node[hwbox, fill=u250red, minimum width=1.55cm] (wl) at (0.8,3.5)
        {$w_{\mathrm{lo}}$\\signed W4};
    \node[hwbox, fill=u250blue, minimum width=3.55cm] (pack) at (4.5,4.5)
        {$P=(w_{\mathrm{hi}}\ll23)+w_{\mathrm{lo}}$\\signed 27 bit};
    \node[hwbox, fill=u250green, minimum width=2.25cm] (act) at (4.5,2.0)
        {$x$ signed A15\\sign-extend 18 bit};
    \node[dspbox] (dsp) at (9.5,4.5)
        {DSP48E2\\$27\times18$ MAC};
    \node[hwbox, fill=u250orange, minimum width=2.65cm] (acc) at (13.5,4.5)
        {packed accumulator\\signed 48 bit};
    \node[hwbox, fill=u250purple, minimum width=4.15cm, minimum height=1.55cm]
        (split) at (18.5,4.5)
        {$y_{\mathrm{lo}}=\mathrm{acc}[22{:}0]$\\
         $y_{\mathrm{hi}}=\mathrm{acc}[45{:}23]+\mathrm{acc}[22]$\\
         hai tích vô hướng signed độc lập};

    \draw[dataline] (wh.east) -| (pack.north);
    \draw[dataline] (wl.east) -| (pack.south);
    \draw[dataline] (pack.east) -- node[above, font=\scriptsize] {operand A} (dsp.west);
    \draw[dataline] (act.east) -| node[below right, font=\scriptsize] {operand B} (dsp.south);
    \draw[dataline] (dsp.east) -- (acc.west);
    \draw[dataline] (acc.east) -- (split.west);
    \node[font=\small, align=center, text=u250navy] at (9.5,2.0)
        {2 phép nhân W4$\times$A15/DSP\\32 lane $\times$ 2 DSP/lane = 64 DSP/PE};

    \draw[black!35, thick] (-0.1,1.0) -- (21.5,1.0);
    \node[font=\bfseries, anchor=west] at (-0.1,0.5)
        {(b) DATAFLOW load--compute--write của linear PE $n$};

    \node[streambox, fill=u250gray] (ddr) at (0.55,-0.5)
        {DDR$n$\\weight W4};
    \node[streambox] (load) at (2.75,-0.5)
        {AXI load\\512b/cycle};
    \node[streambox, fill=u250blue] (fifo) at (5.3,-0.5)
        {BRAM FIFO\\d:512 ($\ge 2$ tile)};
    \node[streambox, fill=u250orange, minimum width=2.0cm] (mac) at (8.5,-0.5)
        {32 lane unroll\\64 DSP (II=1)};
    \node[streambox, fill=u250purple] (partial) at (11.5,-0.5)
        {partial FP32\\128 bit pkt};
    \node[streambox, fill=u250purple] (reduce) at (14.2,-0.5)
        {pair/global\\reduction};
    \node[streambox] (store) at (16.8,-0.5)
        {store local\\output};
    \node[streambox, fill=u250gray] (proj) at (19.4,-0.5)
        {projection\\BRAM};
    \node[streambox, fill=u250green, minimum width=2.8cm, minimum height=0.75cm]
        (as) at (8.5,0.7) {activation A15 + scale\\BRAM/URAM cục bộ};

    \draw[dataline] (ddr) -- (load);
    \draw[dataline] (load) -- (fifo);
    \draw[dataline] (fifo) -- (mac);
    \draw[dataline] (as) -- (mac);
    \draw[dataline] (mac) -- (partial);
    \draw[dataline] (partial) -- (reduce);
    \draw[dataline] (reduce) -- (store);
    \draw[dataline] (store) -- (proj);

    \node[draw=u250navy, thick, rounded corners=2pt, fill=white,
        minimum width=16.8cm, minimum height=0.62cm, align=center,
        font=\small] at (10.0,-2.0)
        {Trạng thái ổn định: \textbf{load tile $n+1$}\quad$\parallel$\quad
         \textbf{compute tile $n$}\quad$\parallel$\quad
         \textbf{reduce/write packet $n-1$}};
\end{tikzpicture}%
}
\caption{Signed W4A15 DSP packing và dataflow load-compute-write của một linear PE.}
\label{fig:w4a15-dataflow}
\end{figure}

\section{Chi Tiết Kiến Trúc Từng Khối Theo Luồng Suy Luận Llama}

Kiến trúc phần cứng được thiết kế tinh gọn để thực thi chính xác luồng suy luận (inference flow) của mô hình Llama-2. Thay vì dùng các module rời rạc cho từng lớp, lõi \texttt{int4\_decoder\_local\_pe} thi hành một vòng lặp gồm 32 lớp (layer), mỗi lớp lặp qua nhiều giai đoạn (stage) tuần tự. 

Mỗi giai đoạn được quyết định bởi một \texttt{mode} tĩnh (Q, K, V, O, Gate, Up, Down, Logits). Cấu trúc \texttt{if-else} trong kernel sẽ tự động bật/tắt các khối xử lý phần cứng phù hợp với \texttt{mode} đó. Dưới đây là phân tích chi tiết từng khối theo đúng chu trình chảy của dữ liệu.

\subsection{1. Khối Chuẩn Hóa RMSNorm và Lượng Tử Hóa (RMS\_STAGE)}

Trước khi đi vào khối Attention hay FFN, tensor đầu vào (residual) phải đi qua RMSNorm. Khối \texttt{int4\_rmsnorm\_quantize\_shards} đảm nhận nhiệm vụ này.

\textbf{Quá trình thực thi:}
\begin{enumerate}
    \item \textbf{Tính tổng bình phương (SumSq):} Mỗi PE tự đọc BRAM \texttt{residual} cục bộ, bình phương các phần tử và tính tổng cục bộ (\texttt{partial\_stream}).
    \item \textbf{Mạng lưới Reduction (Global Sum):} Tổng của PE0+PE1 và PE2+PE3 được cộng chéo giữa các SLR để tạo ra tổng toàn cục, sau đó tính nghịch đảo căn bậc hai ($1/\sqrt{x}$) bằng hàm \texttt{hls::rsqrtf}. Kết quả được phát ngược lại (broadcast) về 4 PE.
    \item \textbf{Chuẩn hóa và Lượng tử hóa INT4:} Mỗi PE dùng giá trị nghịch đảo này nhân với \texttt{residual} và tham số $\gamma$ (được lưu trong URAM \texttt{norm\_cache}). Kết quả FP32 lập tức được lượng tử hóa nhóm (Group Quantization - block size 32) để sinh ra \texttt{activation\_q} (chuẩn INT4) và \texttt{activation\_scale} (FP32).
\end{enumerate}

\subsection{2. Khối Chiếu Tuyến Tính Q, K, V và Lưư BRAM}

Khi \texttt{activation\_q} đã sẵn sàng, hàm \texttt{LINEAR\_STAGE} được gọi 3 lần tuần tự với các \texttt{mode = Q, K, V}.

\textbf{Quá trình thực thi:}
\begin{enumerate}
    \item Động cơ tính toán \texttt{int4\_sharded\_linear\_4pe} nạp trọng số INT4 từ DDR cục bộ.
    \item Thực hiện nhân ma trận (MAC) trên DSP48E2. Partial sums được truyền qua mạng Reduction (chỉ 128-bit đi qua biên SLR).
    \item Kết quả hoàn chỉnh được nạp vào BRAM nháp \texttt{projection}.
    \item Hàm \texttt{int4\_save\_local\_projection} sẽ sao chép kết quả từ \texttt{projection} sang các BRAM chuyên dụng \texttt{q, k, v}.
\end{enumerate}

\subsection{3. Khối Attention (SwiftKV) và Cập Nhật KV Cache}

Khi \texttt{mode = O}, khối \texttt{ATTENTION\_STAGE} (\texttt{int4\_swiftkv\_attention\_pe}) được kích hoạt.

\textbf{Quá trình thực thi:}
\begin{enumerate}
    \item \textbf{RoPE (Rotary Positional Embedding):} Áp dụng biến đổi xoay vị trí cho Q và K bằng các hệ số lượng giác nạp từ \texttt{rope\_lut} (DDR).
    \item \textbf{Cập nhật KV Cache:} K và V mới được ghi thẳng ra bộ nhớ DDR (\texttt{kv\_cache}). Trạng thái Cache của các token trước cũng được nạp song song vào FPGA.
    \item \textbf{Online Softmax Attention:} Phép tính $Q \times K^T$ và $\times V$ được chạy cục bộ trên mỗi SLR. Do kiến trúc Tensor Parallel đã chia head đều cho 4 SLR, bước Attention không cần liên lạc xuyên SLR.
    \item Đầu ra của Attention được lượng tử hóa trực tiếp thành \texttt{activation\_q} mới, chuẩn bị cho lớp Linear O.
\end{enumerate}

\begin{figure}[H]
\centering
\resizebox{\textwidth}{!}{
\begin{tikzpicture}[x=1cm,y=1cm]
    % Memories
    \node[localbox, fill=u250orange] (bram_qkv) at (0, 1.5) {\texttt{BRAM q, k, v}};
    \node[ddrbox, minimum width=3cm] (ddr_rope) at (0, 3) {\texttt{DDR rope\_lut}};
    \node[ddrbox, minimum width=3cm] (ddr_kv) at (0, -1.5) {\texttt{DDR kv\_cache}};
    
    % Attention module detailed
    \node[pebox, fill=u250blue, minimum width=3cm, minimum height=4cm] (attn_mod) at (5.5, 0) {};
    \node[font=\bfseries, text=u250navy, align=center] at (5.5, 1.4) {\texttt{int4\_swiftkv\_}\\\texttt{attention\_pe}};
    
    \node[hwbox, fill=white, minimum width=2.5cm, font=\scriptsize] (rope) at (5.5, 0.4) {RoPE\\Apply};
    \node[hwbox, fill=white, minimum width=2.5cm, font=\scriptsize] (kv_update) at (5.5, -0.6) {KV Cache\\Update};
    \node[hwbox, fill=white, minimum width=2.5cm, font=\scriptsize] (softmax) at (5.5, -1.6) {Online Softmax\\\& Mac};
    
    % Out memories
    \node[localbox, fill=u250orange] (bram_act_q) at (10.5, 0.5) {\texttt{BRAM activation\_q}\\(INT4)};
    \node[localbox, fill=u250orange] (bram_act_s) at (10.5, -0.5) {\texttt{BRAM activation\_scale}\\(FP32)};
    
    % Linear O
    \node[pebox, fill=u250blue, minimum width=2.5cm, align=center] (lin_o) at (15, 0) {\texttt{int4\_sharded\_}\\\texttt{linear\_4pe}\\\texttt{(mode = O)}};
    \node[localbox, fill=u250orange, align=center] (bram_proj) at (19.5, 0) {\texttt{BRAM projection}\\(INT4\_output\_word\_t)};
    \node[pebox, fill=u250green, minimum width=2.5cm] (add) at (19.5, -2.5) {\texttt{int4\_residual\_}\\\texttt{add\_shards}};
    \node[localbox, fill=u250orange] (bram_res) at (15, -2.5) {\texttt{BRAM residual}};
    
    % Lines
    \draw[dataline] (bram_qkv.east) -- node[above, font=\scriptsize] {Q, K, V} (attn_mod.west |- bram_qkv.east);
    \draw[dataline] (ddr_rope.east) -| (rope.north);
    \draw[bothline] (ddr_kv.east) -| (kv_update.south);
    
    \draw[dataline] (attn_mod.east |- bram_act_q.west) -- (bram_act_q.west);
    \draw[dataline] (attn_mod.east |- bram_act_s.west) -- (bram_act_s.west);
    
    \draw[dataline] (bram_act_q.east) -- (lin_o.west |- bram_act_q.east);
    \draw[dataline] (bram_act_s.east) -- (lin_o.west |- bram_act_s.east);
    
    \draw[dataline] (lin_o.east) -- (bram_proj.west);
    \draw[dataline] (bram_proj.south) -- (add.north);
    \draw[bothline] (bram_res.east) -- (add.west);
\end{tikzpicture}
}
\caption{Chi tiết interface và luồng dữ liệu cấp module của khối Attention và Linear O.}
\end{figure}

\subsection{4. Khối Chiếu Đầu Ra (Linear O) và Cộng Residual Đầu Tiên}
Tiếp tục trong \texttt{mode = O}, \texttt{LINEAR\_STAGE} nhận \texttt{activation\_q} (kết quả Attention) và nhân với trọng số $W_O$. Kết quả được xuất ra BRAM \texttt{projection}.
Ngay sau đó, khối \texttt{RESIDUAL\_ADD} được gọi: \texttt{residual = residual + projection}.

\subsection{5. Khối Mạng Truyền Thẳng (FFN: Gate, Up, SwiGLU, Down)}

Nửa sau của layer là khối FFN. Nó bắt đầu bằng một bước RMSNorm thứ hai để sinh ra \texttt{activation\_q}.
\begin{enumerate}
    \item \textbf{Mode GATE:} \texttt{LINEAR\_STAGE} tính toán ma trận Gate, lưu tạm vào BRAM \texttt{projection}. Sau đó hàm save chép nó sang BRAM \texttt{gate} chuyên dụng.
    \item \textbf{Mode UP:} \texttt{LINEAR\_STAGE} tính toán ma trận Up, ghi đè lên BRAM \texttt{projection}. Không cần chép đi đâu nữa.
    \item \textbf{Mode DOWN (Kích hoạt SwiGLU):} Lúc này, khối \texttt{SWIGLU\_STAGE} được gọi. Nó lấy dữ liệu từ BRAM \texttt{gate} và BRAM \texttt{projection} (đang chứa Up) để tính $\text{Gate} \times \sigma(\text{Gate}) \times \text{Up}$. Kết quả lập tức được lượng tử hóa thành \texttt{activation\_q}.
    \item \textbf{Mode DOWN (Chiếu tuyến tính):} \texttt{LINEAR\_STAGE} nhân \texttt{activation\_q} với trọng số $W_{\text{Down}}$, lưu vào BRAM \texttt{projection}.
    \item Cuối cùng, \texttt{RESIDUAL\_ADD} cộng \texttt{projection} vào dòng chảy chính \texttt{residual}.
\end{enumerate}

\begin{figure}[H]
\centering
\resizebox{\textwidth}{!}{
\begin{tikzpicture}[x=1cm,y=1cm]
    \node[localbox, fill=u250orange] (bram_res) at (0, 0) {\texttt{BRAM residual}};
    \node[pebox, minimum width=2.5cm, align=center] (rms) at (4.5, 0) {\texttt{int4\_rmsnorm\_}\\\texttt{quantize\_shards}};
    \node[localbox, fill=u250orange, align=center] (bram_act) at (9.0, 0) {\texttt{BRAM activation\_q/scale}\\(Dùng chung)};
    
    \node[pebox, minimum width=2.5cm, align=center] (gate) at (14.0, 1.5) {\texttt{int4\_sharded\_}\\\texttt{linear\_4pe}\\\texttt{(mode = GATE)}};
    \node[pebox, minimum width=2.5cm, align=center] (up) at (14.0, -1.5) {\texttt{int4\_sharded\_}\\\texttt{linear\_4pe}\\\texttt{(mode = UP)}};
    
    \node[localbox, fill=u250orange, align=center] (bram_gate) at (18.5, 1.5) {\texttt{BRAM gate}};
    \node[localbox, fill=u250orange, align=center] (bram_proj) at (18.5, -1.5) {\texttt{BRAM projection}\\(Chứa giá trị Up)};
    
    \node[pebox, minimum width=2.5cm, align=center] (swiglu) at (23.0, 0) {\texttt{int4\_swiglu\_}\\\texttt{quantize\_shards}};
    
    \draw[dataline] (bram_res) -- (rms);
    \draw[dataline] (rms) -- (bram_act);
    \draw[dataline] (bram_act.east) -- ++(1.5,0) |- (gate.west);
    \draw[dataline] (bram_act.east) -- ++(1.5,0) |- (up.west);
    
    \draw[dataline] (gate) -- (bram_gate);
    \draw[dataline] (up) -- (bram_proj);
    
    \draw[dataline] (bram_gate.east) -| node[above left, font=\scriptsize] {\texttt{gate[i]}} (swiglu.north);
    \draw[dataline] (bram_proj.east) -| node[below left, font=\scriptsize] {\texttt{up[i]}} (swiglu.south);
    
    \node[localbox, fill=u250orange, align=center] (bram_act2) at (23.0, -3.5) {\texttt{BRAM activation\_q/scale}\\(Ghi đè)};
    \node[pebox, minimum width=2.5cm, align=center] (down) at (18.5, -3.5) {\texttt{int4\_sharded\_}\\\texttt{linear\_4pe}\\\texttt{(mode = DOWN)}};
    \node[localbox, fill=u250orange, align=center] (bram_proj2) at (13.0, -3.5) {\texttt{BRAM projection}\\(Ghi đè)};
    \node[pebox, minimum width=2.5cm, align=center, fill=u250green] (add) at (6.5, -3.5) {\texttt{int4\_residual\_}\\\texttt{add\_shards}};
    
    \draw[dataline] (swiglu) -- (bram_act2);
    \draw[dataline] (bram_act2) -- (down);
    \draw[dataline] (down) -- (bram_proj2);
    \draw[dataline] (bram_proj2) -- (add);
    
    \draw[dataline] (bram_res.south) |- node[below right, font=\scriptsize] {\texttt{residual[i]}} (add.west);
    \draw[dataline] (add.north) -- ++(0,1.2) -| node[above, font=\scriptsize] {Ghi đè cập nhật} (bram_res.south east);
\end{tikzpicture}
}
\caption{Chi tiết interface và luồng dữ liệu cấp module của khối FFN.}
\end{figure}

\subsection{6. Logits Head}
Sau 32 layer, \texttt{mode = LOGITS} được gọi. Khối tính thêm một lớp RMSNorm cuối cùng, chạy Linear Logits và lưu mảng xác suất token ra bộ nhớ DDR thông qua \texttt{logits\_pe}.

\section{Giao tiếp Hệ thống, Điều khiển và Ràng buộc Vật lý (Floorplan)}

\subsection{1. Giao tiếp giữa CPU (Host) và FPGA Kernel}
CPU giao tiếp với kernel FPGA thông qua một cổng điều khiển chuẩn \textbf{AXI4-Lite} (\texttt{s\_axilite}). 
Thông qua cổng này, CPU truyền các tham số vô hướng (scalar) và địa chỉ bộ nhớ DDR:
\begin{itemize}
    \item \textbf{Tham số suy luận:} \texttt{position} (vị trí token hiện tại cần giải mã).
    \item \textbf{Địa chỉ vùng nhớ (Pointers):} Địa chỉ vật lý của các băng bộ nhớ \texttt{model\_bank}, \texttt{rope\_lut}, \texttt{kv\_cache}, \texttt{residual} và \texttt{logits}.
    \item \textbf{Tín hiệu Handshake:} Thanh ghi \texttt{ap\_ctrl\_hs} để CPU ra lệnh \textbf{Start} cho kernel và chờ tín hiệu \textbf{Done} khi hoàn thành 32 layer.
\end{itemize}
Khi Kernel đang chạy, nó hoạt động hoàn toàn tự trị (đọc lệnh, giải mã toàn bộ mạng) mà không cần CPU can thiệp cho đến khi sinh ra được token logit cuối cùng.

\subsection{2. Mạng lưới Điều khiển Nội bộ và Giao tiếp Liên SLR}
Nếu sử dụng một Máy trạng thái trung tâm (Central FSM) để phát tín hiệu điều khiển (ví dụ: \texttt{start\_layer\_1}, \texttt{mode=Q}) tới toàn bộ 4 PE trên 4 SLR, tín hiệu sẽ chịu hiệu ứng \textbf{Fan-out toàn cục} khổng lồ, làm sập xung nhịp (timing violation).

Thay vào đó, kiến trúc này sử dụng \textbf{Mạng chuỗi Token (Token Control Chain)}:
\begin{itemize}
    \item \textbf{Lệnh nối tiếp (Relay):} Lệnh vị trí (\texttt{position}) hoặc tín hiệu bắt đầu khối (\texttt{block\_token}) được đưa vào SLR0. SLR0 đọc lệnh, đồng thời chuyển tiếp (relay) lệnh đó sang SLR1 qua một FIFO nhỏ (sâu 2 phần tử). SLR1 lại chuyển tiếp cho SLR2, v.v.
    \item \textbf{Token Hoàn thành (Completion):} Khi một PE tính xong, nó nhả ra một \texttt{completion\_token} (1-bit). Token của PE0 và PE1 được gộp (Join) tại SLR1. Token của PE2 và PE3 được gộp tại SLR2. Cuối cùng, hai tín hiệu gộp này lại hội tụ về SLR1 để xác nhận toàn bộ 4 PE đã xong.
    \item \textbf{Lợi ích:} Nhờ cơ chế chuyền tay (như chạy tiếp sức), không có bất kỳ dây điều khiển nào chạy xuyên qua nhiều hơn 1 biên giới SLR. Các controller của các SLR hoàn toàn tách rời và tự kích hoạt khi có token đẩy vào FIFO của nó.
\end{itemize}

\subsection{3. Chiến lược Floorplan Hạn chế SLL và Tối ưu Timing (Fmax > 270 MHz)}
Kiến trúc FPGA ghép mảnh (Multi-SLR) giao tiếp với nhau qua các đường Super Long Lines (SLL). Nếu các logic bị bộ tổng hợp (Synthesis) vứt rải rác lộn xộn qua các SLR, số lượng SLL sẽ bùng nổ, gây nghẽn định tuyến (routing congestion) và tụt Fmax thảm hại.

Dự án này giải quyết bài toán Fmax $\geq$ 300MHz bằng cách áp dụng \textbf{Hard Floorplan} thông qua các file cấu hình TCL (\texttt{timing\_300mhz\_domains.tcl} và \texttt{link\_300mhz.cfg}):
\begin{enumerate}
    \item \textbf{Giới nghiêm (Containment) bằng Pblock:} Script TCL sẽ chủ động duyệt cây phân cấp (hierarchy) của file RTL sau khi Vitis HLS sinh ra. Nó gom tất cả logic, BRAM, URAM, DSP thuộc về \texttt{PE0} và ép cứng (assign) vào vùng \texttt{pblock\_dynamic\_SLR0}. Tương tự cho \texttt{PE1} vào SLR1, \texttt{PE2} vào SLR2, \texttt{PE3} vào SLR3.
    \item \textbf{Gắn chặt AXI với DDR vật lý:} Cổng \texttt{m\_axi\_gmem0} được ép kết nối trực tiếp với bộ nhớ DDR[0] nằm ngay trên cùng một SLR0, loại bỏ hoàn toàn việc đường truyền AXI 512-bit bị vắt chéo SLR.
    \item \textbf{Kiểm soát chặt chẽ biên giới SLR:} Các điểm cắt SLL duy nhất được phép tồn tại là:
    \begin{itemize}
        \item Đường truyền Token điều khiển hẹp (12-bit cho position, 1-bit cho completion).
        \item Đường truyền Partial Sum (128-bit) của mạng Reduction. 
    \end{itemize}
    Tất cả các đường này đều được ép phải đi qua thanh ghi chốt (Pipeline Registers / FIFO \texttt{impl=bram}) ngay tại biên giới SLR.
\end{enumerate}
Nhờ việc định hướng trước mọi đường đi (Data-driven KPN) và ép cứng vị trí linh kiện (Strict Floorplanning), bộ định tuyến (Router) của Vivado không phải đoán mò. SLL được giảm thiểu tối đa, giúp luồng thiết kế dễ dàng vượt mốc tần số \textbf{300 MHz} mà không vi phạm timing.

\section{Kết luận}
Kiến trúc phần cứng được hiện thực hóa mang tính tối ưu cao cho luồng xử lý suy luận LLM theo phương thức memory-bound. Việc tái sử dụng lõi tính toán, kết hợp lượng tử hóa INT4 và đóng gói hiệu năng vào DSP48 biến Alveo U250 trở thành một cỗ máy gia tốc mạnh mẽ cho mô hình Llama-2 7B.

\end{document}
